\chapter{A Difference Appears}

\section{The first fact}

Before a symbol can be measured it must first be told apart from what it is not, and that telling-apart is the whole
of the machine's ground floor. Nothing here is counted, encoded, or priced. There is only a \emph{this} that is not a
\emph{that} --- the smallest possible act, prior to number, on which every later act is built. This section builds that
act and only that act, and it is worth building slowly, because everything the machine will ever compute is this one
decision, iterated.

In this volume the first fact is read as a piece of a formal system, and the reading pays off at once. A difference
here is not an event that happens and is then recorded; it is a \emph{decidable predicate}. To tell two symbols apart
is to run a finite procedure that halts and reports \emph{not-equal} --- a decision the machine performs, not a fact
it files away. The distinction is not asserted and then checked against evidence later; the running \emph{is} the
evidence. A difference the machine could not decide in finite time would be no difference the machine could use, and so
at the ground floor a difference is precisely what a halting procedure certifies: two marks, and a computation that
tells them apart.

This is the coarsest object information theory knows, and it is the one everything else is built on. A symbol carries
meaning, Shannon~\cite{shannon1948} showed, exactly to the degree it could have been another --- a mark on a channel
informs you precisely because the channel might have carried a different mark in its place. Strip that idea to its
floor and you reach this machine's first fact: not yet a quantity of information, not yet a rate or an entropy, but the
bare precondition for any of them --- that two symbols can be distinguished at all. Where Shannon counts how much a
symbol could have been otherwise, the ground floor asks only \emph{whether} it could: whether the alphabet holds a
second mark this one is not. The bit, the entropy, the channel capacity are all later constructions. The floor beneath
them is a single decided difference.

And here is the first of the book's two working verbs, introduced where it is native. To \emph{tange} is to select ---
to pick a \emph{this} out of an alphabet by the decidable characteristic that parts it from its \emph{that}. Telling
two symbols apart is not yet counting them; it is the pure act of decidable selection, and that act is a tange. The
verb earns its keep exactly here, because the machine's first faculty is nothing but a decision procedure that, handed
two marks, halts and selects: \emph{these differ}, or \emph{they do not}. Its partner verb --- to \emph{funge}, to
pool a mark with all the marks a decision cannot separate from it --- is not yet available, and cannot be, for you
cannot pool by sameness what you have not first tanged apart. That is the next chapter's work.

This lands the reader on the one fact the ground floor exists to establish: \emph{resolution is logically upstream of
measurement}. You must be able to decide a symbol apart before you can have a number of symbols; distinguishability
comes before quantity. It is the same point Nyquist~\cite{nyquist1928} and later Shannon~\cite{shannon1949} made about
signals --- that you cannot read past a resolution your channel cannot resolve --- pressed here one level lower, to
symbols themselves: you cannot read a scale whose marks you cannot first decide into distinct symbols. The mark must be
a \emph{this-not-that} before it can be a \emph{how-many}.

It is worth marking how much of the book is already latent in this one decision. Everything the machine will later call
a number, a code, a cost, a bound is built by iterating exactly this: decide a difference, let the decision stand as
its own record, and pay only for the running. The naming of the electron eight chapters on, and the honest bracket the
machine hands back at the very end, are this same act matured --- not new machinery, but the first decided difference
taken around the loop enough times to be read. And notice what has \emph{not} appeared: no number, no count, no
magnitude. The first fact is pre-numeric by construction, which makes this the cleanest place to adopt the discipline
the whole instrument keeps --- it will not reach for a quantity it cannot decide, and the reason is visible already
here, at the floor, where there is not yet even a number to reach for.

\section{The symbol as a label}

The first fact was a decided difference; the natural next question is where the machine \emph{stores} it. Having
decided two marks apart, where does it write the distinction down, so that it can find it again? The answer is the
hinge of this chapter, and it is a small shock: nowhere. The machine writes the difference nowhere, because the
difference is already its own record. The distinction and the label of the distinction are not two things joined by an
act of naming; they are one object. And because they are one object, the label costs nothing to keep. In this machine,
difference is free.

To see why, ask what the decision of the last section actually \emph{produced}. In the ordinary account there are two
separate things with a copy between them: a difference on one side, a name for the difference filed somewhere on the
other, and a step that carries the first into the second. Here that middle step is missing, and its absence is the
whole content of the section. What the machine reads is a \emph{codeword} --- the position of the symbol in the
alphabet, the very characteristic the decision selected by. Deciding the two marks apart already \emph{produced} that
position; there is no second act in which the machine turns around and records what it found. The tange did not select
a difference and then go write its name; the selection \emph{was} the name.

This is why the claim ``the difference is its own label'' is not a play on words. In this register a label just
\emph{is} a codeword --- a position in the alphabet, which mark relative to which --- and the tange delivered exactly
such a position. To decide a \emph{this} against a \emph{that} is already to hold the codeword that parts them; and the
codeword is the label. Naming adds nothing, because the decision already named. The characteristic the machine selected
by and the record it walks away with are the same object, and the reader who was holding them apart as two may let them
fall together into one. A good code, Shannon and his successors would say, wastes no symbol pinning a name onto a
distinction the distinction already carries; this machine's code is that thrift taken to its floor.

The consequence is that the distinction requires no storage of its own --- but the claim must be scoped exactly,
because scoped loosely it is false. The machine does not reserve a cell \emph{for the difference}; it reads the
distinction straight off the one index it must track anyway to keep its own descriptions coherent, the position in its
own alphabet. It does not follow that the machine reserves \emph{nothing}. The construction carries a single standing
datum, one shared constant against which decisions are taken, and it is honest about it: that one fixture lives on the
data page, not spent afresh on each telling-apart. Read the carve-out as carefully as the claim. What is free is the
\emph{difference} --- nothing is set aside per distinction, because the codeword lives at a position the machine
already keeps --- while the one thing allocated, the shared constant, is a single fixture, not a fresh cell per
decision. ``Free'' here does not mean fast; that would be a claim about a clock, and this is not one. It means
\emph{nothing added}: the difference is decided where it stands, not stored to be looked up.

Underneath all of this is an idea far older than any computer, and the machine takes it at its word. A symbol means
nothing in isolation and means only by its contrast with the symbols it is not --- Saussure~\cite{saussure1916} said it
of words, that value is difference, and the theory of communication later made it quantitative. This machine builds on
the old idea literally: its symbols \emph{are} their differences, a mark defined by its inequality to every other mark.
The label is the contrast itself, not a name pinned on after the contrast was noticed. The semiotics is illumination
here, texture around the fact; what the machine actually carries is the decided inequality, and nothing of Saussure's
beyond it.

So take stock of what the ground floor has handed you: not a measurement, and not yet a number --- a \emph{position}.
Which mark, relative to which. That is all of it, and it is exactly enough, because a symbol decided apart from other
symbols is the seed of everything the machine will grow. And the label is about to earn its keep a second way: in the
very next chapter the machine will \emph{funge} --- pool marks that share a codeword into a single bin --- and the size
of that bin will be its first number. The codeword is precisely what the count will pool by; but that is the next
chapter's work, and it performs no count here. What this section settles is only that the step onto the count carries
no hidden storage to be paid for later. Everything the machine will grow is this one act iterated: decide a difference,
let it be its own codeword, and pay only for the running.

\section{What a difference is not}

The ground floor is nearly built, and the last thing to do on it is the most easily skipped: to say what the first
fact is \emph{not}. A difference is not yet a number. It is not a code length, not a rate, not a distance --- the
smallest possible act, and no more. The temptation, having built something, is to read more into it than it holds; the
whole discipline of this book is the refusal of that temptation, and the refusal is easiest to learn here, where there
is least to inflate.

The relation the machine actually has is inequality, and inequality alone. To ask whether two marks differ is to ask
whether one is unequal to the other, and the decision returns \emph{same} or \emph{not-same} and halts. Inequality is
the coarsest relation two symbols can stand in --- a bare yes-or-no that carries nothing else along with it. It is
worth pausing on how much that ``nothing else'' rules out, because each thing it rules out is a structure the reader
may be tempted to assume is already present. There are three, and each is a refusal.

The first refusal: a difference is \emph{not an order}. To be a number is, at the least, to sit in a chain --- to be
comparable, to have a \emph{less-than} and a \emph{greater-than}. The bare distinction exposes no such thing. Its one
relation is symmetric and blind to direction: it can report that two marks are not the same, but not which comes first,
or is larger, or sits higher. Order does not live on the distinction at all; it is built higher up, on the numeric
objects the later chapters construct --- never as a feature of the raw difference. A reader who searches the ground
floor for a \emph{which-is-bigger} will not find one, and not because it is hidden but because the question cannot yet
be posed. At this level, ``which is bigger?'' does not so much as parse.

The second refusal: a difference is \emph{not a magnitude}. A measurement answers \emph{how much}; a difference
answers only \emph{whether}. There is no unit here, no scale, no metric --- nothing that would let the machine call two
marks far apart or close, only that they are two. The datum is nominal in the plainest sense: same or different, and no
finer. One might recall Frege's~\cite{frege1884} remark that a number belongs to a concept and not to a bare object ---
but only in passing, for the machine's refusal here is not a philosopher's; it is structural, a matter of which fields
the ground floor does and does not carry, and it carries no metric at all.

The third refusal is the one this book is built to make precise, and it is cleanest stated in the two working verbs. A
difference is a pure \emph{tange} --- a decidable selection, a \emph{this} decided apart from a \emph{that}. A count is
something else entirely: it is the size of a \emph{funge}, the cardinality of a bin of like marks. And you cannot take
the size of a bin before you have a bin, nor pool marks by their likeness before you have first tanged them apart. So
``not yet a number'' is not a shortfall the machine will outgrow by trying harder; it is an \emph{ordering}, as strict
as any in the construction. The tange comes first; the funge comes after; and number lives entirely on the funge side
of that line. A number is a \emph{later construction} --- something the next chapter builds by pooling marks the ground
floor has already tanged apart --- not a reading available here.

That, in the end, is the shape of the whole discipline, seen in miniature at the smallest possible scale. The machine's
refusal to call a position a quantity is not timidity, and it is not incompleteness; it is correctness. It holds exactly
one thing --- a distinction, a codeword, which mark relative to which --- and it declines to inflate that one thing into
a number it has not yet built. This is the same refusal, grown up, that the instrument will make at the very end, when
it declines to call its honest bracket the real value the bracket contains. In both places the machine reports exactly
what it has decided and not a hair more. The ground floor and the final floor keep the same discipline; what stands
between them is only the long climb --- and the climb begins in the next chapter, when the machine takes the differences
it has tanged apart and, for the first time, begins to funge them into a count.
